% Teaching manuscript. Edit this TeX source; Markdown is compatibility navigation only.

\StudyChapter{CS5222-Tutorial04}{Tutorial 4 · File Distribution, BitTorrent and DASH}
\StudyTarget{CS5222-Tutorial04__tutorial-4-file-distribution-bittorrent-and-dash}



\StudyHeading{1}{上传能力与文件分发时间}
\StudyTarget{CS5222-Tutorial04___1}


\StudyRef{CS5222-Chapter02}{root}{Chapter 2} · \StudyRepoRef{CS5222/course-notes/README.md}{}{课程目录} · \StudyRef{CS5222-Tutorial05}{root}{下一份：Tutorial 5}


前两题先找分发时间的下界，再说明怎样分配上传速率才能达到它；后两题分别检查BitTorrent的服务机会和DASH的音视频组合条件。

依据原题2页和教师解答v2共15页，按Problem1–4。本Tutorial配套Chapter2应用层。题目单位用十进制Gbit/Mbit/kbit。


\textbf{按需基础：} \StudyRef{Foundations-NetworkBasics}{units}{单位换算}、\StudyRef{Foundations-NetworkBasics}{timeline}{速率与完成时间}。先会用“任务量\ensuremath{\div}每秒处理量”算时间。


对应讲解：\StudyRef{CS5222-Chapter02}{file-distribution}{分发下界}解释三类容量约束；\StudyRef{CS5222-Chapter02}{bittorrent}{BitTorrent请求与供给}区分稀缺块选择和上传机会；\StudyRef{CS5222-Chapter02}{video-cdn}{DASH与CDN}把表示、时间块、服务节点及播放缓冲接起来。


\begingroup
\small\setlength{\tabcolsep}{3pt}
\renewcommand{\arraystretch}{1.22}
\begin{longtable}{@{}>{\raggedright\arraybackslash}p{\dimexpr0.191048\linewidth-1.50000\tabcolsep\relax}>{\raggedright\arraybackslash}p{\dimexpr0.153992\linewidth-1.50000\tabcolsep\relax}>{\raggedright\arraybackslash}p{\dimexpr0.210151\linewidth-1.50000\tabcolsep\relax}>{\raggedright\arraybackslash}p{\dimexpr0.444808\linewidth-1.50000\tabcolsep\relax}@{}}
\toprule
题目 & 优先级 & 难度与卡点 & 掌握要求／依据 \\
\midrule\endfirsthead
题目 & 优先级 & 难度与卡点 & 掌握要求／依据 \\
\midrule\endhead
\bottomrule\endfoot
Q1 分配方案 & 核心必会 & 中，下界还需构造达到它的方案 & 按两种瓶颈给实际速率；Tutorial 4教师解答 v2，第1–5张幻灯片 \\
Q2 时间计算 & 核心必会 & 中，总供给与最慢下载 & 统一单位、比较每个约束；Tutorial 4教师解答 v2，第6–8张幻灯片 \\
Q3 free-riding & 常规掌握 & 中，可能与保证不同 & 解释optimistic unchoking条件；Tutorial 4教师解答 v2，第9–11张幻灯片 \\
Q4 DASH & 核心必会 & 中，配对假设 & 区分N、N\textsuperscript{2}与2N；Tutorial 4教师解答 v2，第12–15张幻灯片 \\
\end{longtable}
\endgroup







\StudyHeading{1}{Problem 1. Construct a distribution scheme}
\StudyTarget{CS5222-Tutorial04__q1}
\StudyTarget{CS5222-Tutorial04__problem-1-construct-a-distribution-scheme}


\textbf{原条件 / Conditions：} 服务器初始持有F bit文件，分给N个客户端，服务器总上传率\(u_s\)，客户端下载率\(d_i\)，\(d_{min}=\min_i d_i\)。流体模型允许同时以不同速率发送，总发送率不超过\(u_s\)，忽略其他瓶颈。 / The server simultaneously allocates divisible traffic to N clients, subject to its total upload and each client's download limit.


服务器必须上传NF bit，所以时间至少NF/\(u_s\)；最慢客户端必须收F bit，所以至少F/\(d_{min}\)。但“至少这么久”还不是“确实做得到”，题目要求给方案。


\textbf{a. \(u_s/N\lt d_{min}\)：} 平均给每个客户端\(u_s/N\)。总速率正好\(u_s\)，每人都不超过其下载上限；每人同时在\(F/(u_s/N)=NF/u_s\)完成。瓶颈是服务器供给。


\textbf{b. \(u_s/N>d_{min}\)：} 每人按\(d_{min}\)发送。总量\(Nd_{min}\lt u_s\)，所有客户端都收得下；同在\(F/d_{min}\)完成。快客户端没有跑满也不妨碍整体最优，因为慢客户端仍决定结束时刻。


合起来，每人分配\(\min(u_s/N,d_{min})\)，达到两项下界的最大值。等号边界两种表达相同。Tutorial 4教师解答 v2，第2张幻灯片/Tutorial 4教师解答 v2，第4张幻灯片的反证是在同一流体约束下说明方案，不是现实TCP性能保证。


\textbf{English answer:} Allocate min(us/N,dmin) to every client. The aggregate rate is feasible and every client completes in max(NF/us,F/dmin).


\textbf{独立题 / Transfer：} F=100Mbit、N=4、\(u_s\)=20Mbps、每人下载8Mbps，按本节流体模型给出可行分配与完成时间。 / For F=100 Mbit, N=4, server upload 20 Mbps and each client download 8 Mbps, give a feasible allocation and completion time.



\begin{StudySolution}{答案 / Answer}


每人5Mbps，20秒。最慢下载约束12.5秒不是最终答案，要与总上传约束20秒取最大。 / 5 Mbps each, finishing in 20 seconds.


\end{StudySolution}






\StudyHeading{1}{Problem 2. 15 Gbit分给100人}
\StudyTarget{CS5222-Tutorial04__q2}
\StudyTarget{CS5222-Tutorial04__problem-2-15-gbit100}


\textbf{完整条件 / Conditions：} F=15Gbit，N=100，服务器30Mbps，每人下载2Mbps、上传700kbps=0.7Mbps。问client-server与P2P理想最短分发时间。 / A 15-Gbit file is distributed to 100 peers. Server upload is 30 Mbps; each peer downloads at 2 Mbps and uploads at 700 kbps. Using decimal units and the ideal model, compare client-server and P2P completion times.


先把文件写成15000Mbit。Client-server两项为100\ensuremath{\times}15000/30=50000秒，以及15000/2=7500秒，因此选\textbf{50000秒，C}。


P2P仍需要服务器发出至少一份文件，至少500秒；每人下载至少7500秒；全体上传总能力30+100\ensuremath{\times}0.7=100Mbps，供给NF需要15000秒。三项取最大，得\textbf{15000秒，B}。


\StudyImage{CS5222/course-notes/assets/oct-distribution.png}
\StudyCaption{分发时间的各项容量下界；柱子的单位是秒，比较最大者。}


\begin{StudyBlock}{读图与条件}{取最大下界，并说明可达条件}
图中每根柱都是一个必要时间约束，最大者决定该理想模型下的瓶颈。不是把三根柱相加，因为不同资源可同时工作。原教师页的下界符号视觉上写成“>”，最短时间公式应理解为“至少”，流体理想条件下可以达到；严格大于会排除等号而与Q1的构造冲突。
\end{StudyBlock}


\textbf{English answer:} Client-server: max(50000,7500)=50000 s. P2P: max(500,7500,15000)=15000 s. Upload assistance reduces the aggregate bottleneck but does not remove the other constraints.


\textbf{变式 / Transfer：} F=15 Gbit、N=10、服务器上传30 Mbps、每人下载2 Mbps且上传0.7 Mbps。按本节理想模型，P2P会比client-server快吗？ / For F=15 Gbit and ten clients, use server upload 30 Mbps, each client download 2 Mbps and upload 0.7 Mbps. In the ideal model, compare P2P and client-server completion-time bounds.



\begin{StudySolution}{答案 / Answer}


client-server=max(5000,7500)=7500秒；P2P=max(500,7500,150000/37)=7500秒。此时每人的2Mbps下载是瓶颈，增加上传互助不会突破它。 / Both are 7500 s, limited by the slowest download rate.


\end{StudySolution}






\StudyHeading{1}{Problem 3. 不上传的人，是否仍可能下载完成？}
\StudyTarget{CS5222-Tutorial04__q3}
\StudyTarget{CS5222-Tutorial04__problem-3}


\textbf{原题 / Task：} Bob加入BitTorrent但不上传；a能否完整下载？b多台不同IP电脑是否可能让这种free-riding更高效？ / Can a non-uploading participant finish, and can multiple hosts improve that possibility?


Tutorial 4教师解答 v2，第9–11张幻灯片答两个问题均\textbf{可能，B}。除了优先服务回馈较好的peer，课堂BitTorrent模型还有optimistic unchoking：周期性试着给某个新peer机会。如果有足够peer停留足够久，Bob可能逐渐收齐；参与者过早离开则未必能完成。


多主机可各自获得机会并收集不同片段后合并，教师称为一种Sybil式行为。这是机制分析；实际速率还受peer数量、停留时间、网络资源等影响，不能把主机数量乘上去就断言下载速度同比增加。


\textbf{English answer:} Both claims are possible under the stated classroom assumptions. Optimistic unchoking can supply a free-rider; multiple identities may collect complementary chunks. Neither completion nor proportional speedup is guaranteed.


\textbf{自测 / Check：} 如果所有peer在Bob收齐之前离开，前面的结论是否矛盾？ / Does early peer departure contradict the possibility claim?



\begin{StudySolution}{答案 / Answer}


不矛盾，原答案有停留时间和供给条件。“可能”不等于“必然”。 / No; the sufficient opportunity may not exist in that swarm.


\end{StudySolution}






\StudyHeading{1}{Problem 4. DASH到底存N份，还是N\textsuperscript{2}份？}
\StudyTarget{CS5222-Tutorial04__q4}
\StudyTarget{CS5222-Tutorial04__problem-4-dashnn2}


原题有N个视频质量版本和N个音频质量版本。开头说任意选择音视频，a的括号又明确要求按质量一一配对；教师解答v2第14张幻灯片按\textbf{一一配对}作答，因此本题a为\textbf{N，A}。b分离保存音频和视频，需\textbf{2N，B}，客户端负责同步。


以N=3看清：固定高配高、中配中、低配低，预合并文件只有3份。若真正允许任意组合，要把3\ensuremath{\times}3种配对都提前合并，则是9份。若分轨，保存3视频+3音频共6份，播放时组合即可。这三种情境不是互相推翻，差别在配对约束。


\textbf{English answer:} Under part a's explicit one-to-one matching assumption, N combined files are needed. Separate audio and video require 2N files. N\textsuperscript{2} applies only when every pair is precombined. These counts concern representations, ignoring segment multiplication.


\textbf{独立题 / Transfer：} 视频4种、音频2种，全部预组合与分轨分别几份？ / With four video versions and two audio versions, how many files are needed when every pair is precombined, and when tracks are stored separately?



\begin{StudySolution}{答案 / Answer}


8与6。不是4\textsuperscript{2}，因为两边版本数量不同。 / Eight combined files or six separate tracks.


\end{StudySolution}



\StudyHeading{1}{复习与来源索引}
\StudyTarget{CS5222-Tutorial04___2}


原题p1含Q1–3，p2含Q4。Tutorial 4教师解答 v2，第1–5张幻灯片对应Q1，Tutorial 4教师解答 v2，第6–8张幻灯片对应Q2，Tutorial 4教师解答 v2，第9–11张幻灯片对应Q3，Tutorial 4教师解答 v2，第12–15张幻灯片对应Q4。原题数值、单位及括号条件均保留；图为按原数值生成的教学图。


复习沿用\href{https://crazyshout.github.io/micro-course/cards.html\#CS5222-N234}{N234}、\href{https://crazyshout.github.io/micro-course/cards.html\#CS5222-N235}{N235}、\href{https://crazyshout.github.io/micro-course/cards.html\#CS5222-N236}{N236}。Q3可结合optimistic unchoking，复述为什么答案是“可能”。
